%
% invtan.tex -- Umkehrfunktion des Tangens
%
% (c) 2021 Prof Dr Andreas Müller, OST Ostschweizer Fachhochschule
%
\documentclass[tikz]{standalone}
\usepackage{amsmath}
\usepackage{times}
\usepackage{txfonts}
\usepackage{pgfplots}
\usepackage{csvsimple}
\usetikzlibrary{arrows,intersections,math}
\definecolor{darkred}{rgb}{0.8,0,0}
\begin{document}
\def\skala{1}
\def\R{2.5}
\def\w{39}
\def\r{1.2}
\begin{tikzpicture}[>=latex,thick,scale=\skala]

\draw[color=gray!50] (0,0) circle[radius=\R];

\fill[color=gray!20] (0,0) -- (-\w:\r) arc(-\w:\w:\r) -- cycle;
\node[color=blue] at ({0.5*\w}:{0.7*\r}) {$x$};
\node[color=blue] at ({-0.5*\w}:{0.7*\r}) {$x$};

\draw[->] (-2.7,0) -- (3,0) coordinate[label={$\operatorname{Re}$}];
\draw[->] (0,-2.7) -- (0,3) coordinate[label={right:$\operatorname{Im}$}];

\draw[color=blue] (0,0) -- (\R,{\R*tan(\w)});
\draw[color=blue] (0,0) -- (\R,{-\R*tan(\w)});

\node[color=darkred] at (\R,{0.5*\R*tan(\w)}) [right] {$iy$};
\node[color=darkred] at (\R,{-0.5*\R*tan(\w)}) [right] {$-iy$};
\node[color=darkred] at ({0.7*\R},0) [above] {$1$};

\draw[color=darkred,line width=1.2pt] (0,0) -- (\R,0);
\draw[color=darkred,line width=1.2pt] (\R,{-\R*tan(\w)}) -- (\R,{\R*tan(\w)});

\fill[color=blue] (\R,{\R*tan(\w)}) circle[radius=0.05];
\fill[color=blue] (\R,{-\R*tan(\w)}) circle[radius=0.05];
\fill[color=darkred] (0,0) circle[radius=0.05];

\node[color=darkred] at (\R,{\R*tan(\w)}) [above] {$1+iy$};
\node[color=darkred] at (\R,{-\R*tan(\w)}) [below] {$1-iy$};

%\node[color=blue] at (\R,{\R*tan(\w)}) [above left] {$e^{ix}$};
%\node[color=blue] at (\R,{-\R*tan(\w)}) [below left] {$e^{ix}$};

\end{tikzpicture}
\end{document}

